\pagestyle{fancy}



\phantomsection
\addcontentsline{toc}{section}{RESUMEN}
\section*{Resumen}


\noindent
El resumen permite identificar la esencia del escrito, mencionando brevemente el objetivo y la metodología, así como los resultados y las conclusiones (mínimo 150, máximo 250 palabras).

% Las palabras clave son los términos, materias y terminología que hacen posible describir y recuperar un documento en una disciplina específica. Pregúntese, por ejemplo: ¿con qué palabras puede un usuario de Internet recuperar mi documento? ¿Cuáles son los términos con los que mis colegas abordan esta temática? 3-7 palabras clave.
\vspace{1cm}
\palabrasClave{Artículo científico, artículo de revisión, investigación, estilos de citación}

% Miembros de la Facultad de Ciencias Económicas incluyen Clasificación JEL de la American Economic Association, disponible en https://bit.ly/3bdELGL
% Ejemplo de formato después de palabras clave:
% Clasificación JEL: Q01, Q13, Q16



% ==================================================



\newpage
\phantomsection
\addcontentsline{toc}{section}{ABSTRACT} 
\section*{Abstract}


\noindent
El abstract es el mismo resumen pero en idioma inglés. Conserva la misma extensión o aproximada, es decir, mínimo 150 y máximo 250 palabras.


% Las keywords son las mismas palabras clave pero en inglés. Dependiendo de la disciplina, indague en bases de datos y plataformas cuáles son las traducciones al inglés de sus palabras clave, tales como IEEE Taxonomy (ingenierías IEEE http://goo.gl/nf29xl) o en diccionarios y tesauros especializados.
\vspace{1cm}
\keywords{Scientific article, review article, research, citation styles}

% Miembros de la Facultad de Ciencias Económicas incluyen Clasificación JEL de la American Economic Association, disponible en https://bit.ly/3bdELGL
% Ejemplo de formato después de keywords:
% JEL Classification: Q01, Q13, Q16